% !PRINTED 27
\weberpage{27}
%<<CONT
und daraus ist nach (12) und (14) abzuleiten:
\[
\begin{aligned}
\Phi' &= M\Phi + N\Psi,\\
\Psi' &= P\Phi + Q\Psi,
\end{aligned}
\]
worin die Koeffizienten $M$, $N$, $P$, $Q$ noch zu bestimmen sind.
Wir können dies durch Rechnung ausführen. Ohne Rechnung ergibt sich
das Resultat auf folgendem Wege:

Die $M$, $N$ sind ganze Funktionen zweiten Grades von $p$, $q$ und
dritten Grades von $m$, $n$. Nimmt man $p:q = m:n$, also
$r = mq - np = 0$, so gehen $\Phi'$ und $\dfrac{\partial \Phi'}{\partial p}$
nach (18) in lineare Verbindungen von $\varphi$ und $\psi$ über und
verschwinden also. Daraus folgt, daß $M$ und $N$ durch $r^2$ teilbar
sind, und die Quotienten sind lineare Funktionen von $m$ und $n$. Da
aber für die beiden Fälle $\begin{pmatrix} m & n\\ p & q \end{pmatrix} = \begin{pmatrix} 1 & 0\\ 0 & 1 \end{pmatrix},\ \begin{pmatrix} 0 & 1\\ 1 & 0 \end{pmatrix}$
$\Phi'$ in $\Phi$ und in $\Psi$ übergehen muß, und da man dieselbe
Betrachtung auf $\Psi'$ anwenden kann, so folgt:
\weq{19}{\begin{aligned}\Phi' &= r^2(m\Phi+n\Psi),\\ \Psi' &= r^2(p\Phi+q\Psi).\end{aligned}}
Es werden also (von dem Faktor $r^2$ abgesehen) die Funktionen
$\Phi$, $\Psi$ mit den $\varphi$, $\psi$ \gesperrt{kongredient} transformiert.
(Über die strenge Begründung dieser Schlüsse sehe man Bd. I, §\,20.)

Aus diesen Ergebnissen können wir auf einfache Weise, immer unter der
Voraussetzung $\varphi = 0$, $\psi = 0$, die $y_i^2$ durch $\Phi$ und
$\Psi$ ausdrücken, also die linearen Gleichungen (12), (14) auflösen.
Wenn wir nämlich die Formeln (19) auf die Substitution
\[
\begin{pmatrix} m, & n\\ p, & q \end{pmatrix} = \begin{pmatrix} \beta_1, & -\alpha_1\\ 0, & 1 \end{pmatrix}
\]
anwenden, so ist $r = \beta_1$, die Funktion $\psi'$ bleibt
ungeändert $= \psi$ und $\varphi'$ geht aus $\varphi$ hervor, wenn
\[
\alpha_1,\quad \alpha_2,\quad \alpha_3,\quad \alpha_4
\]
durch
\[
0,\quad (\alpha_2\beta_1),\quad (\alpha_3\beta_1),\quad (\alpha_4\beta_1)
\]
ersetzt werden. Es wird also nach (14)
\[
\Phi' = \beta_1^2(\alpha_2\beta_1)(\alpha_3\beta_1)(\alpha_4\beta_1)\,y_1^2
\]
und folglich nach (19):
\weq{20}{\begin{aligned}
(\alpha_2\beta_1)(\alpha_3\beta_1)(\alpha_4\beta_1)\,y_1^2 &= \alpha_1\Phi - \beta_1\Psi,\\
(\alpha_1\beta_2)(\alpha_3\beta_2)(\alpha_4\beta_2)\,y_2^2 &= \alpha_2\Phi - \beta_2\Psi,\\
(\alpha_1\beta_3)(\alpha_2\beta_3)(\alpha_4\beta_3)\,y_3^2 &= \alpha_3\Phi - \beta_3\Psi,\\
(\alpha_1\beta_4)(\alpha_2\beta_4)(\alpha_3\beta_4)\,y_4^2 &= \alpha_4\Phi - \beta_4\Psi.
\end{aligned}}
